\chapter{Conclusion}
\label{ch:conclusion}

\sysname{} answers the question posed in Chapter~\ref{ch:introduction}: a rate limiter can
give a bounded, known guarantee without a coordinator on the request path, provided clients
accept a probabilistic bound during the brief window after traffic shifts rather than an
exact one at all times. The results in Chapter~\ref{ch:results} show that window closing
within $O(\log n)$ gossip rounds, and Figure~\ref{fig:latency} shows the excess a client
experiences during it falling below the burst the underlying token bucket already permits.
The price is small and the latency win, shown in Table~\ref{tab:throughput}, is not.

\section{Limitations}

The evaluation in Chapter~\ref{ch:results} uses a synthetic workload with one traffic shift
per run; production traffic shifts continuously, and it remains open whether overlapping
shifts compound the excess-admission fraction or whether the estimator's exponential
decay term keeps them independent. The construction also assumes a gossip layer already
exists in the deployment, which is true of the systems we evaluated against but is an
additional operational dependency for a deployment starting from nothing.

\section{Future work}

Three extensions seem tractable. First, the estimator in
Section~\ref{sec:distributed-bucket} assumes uniform node weight; weighting gossip
exchanges by a node's recent traffic share might shorten convergence further under skew.
Second, the token bucket parameters $C$ and $r$ from Section~\ref{sec:token-bucket} are
fixed for the lifetime of a deployment in this thesis; making them adaptive to observed
traffic is a natural next step. Third, the evaluation never asks what an adversarial client
can do to the estimator itself, and a treatment of gossip poisoning would be needed before
\sysname{} could be recommended for a rate limiter facing untrusted clients directly.
